\hrule % Add horizontal line
\vspace{3mm}
\justifying
    \textbf{Digital Advertising Auctions with Asymmetric Consumer Information} \\[3pt]
    With \href{https://www.eui.eu/people?id=rodolfo-salonia}{Rodolfo Alberto Salonia} (EUI) \\[3pt]
    We study how preferential access to consumer data shapes competition and data management strategies in digital display advertising markets. We develop a model in which two demand-side platforms (DSPs) compete to purchase a single advertising slot through a first-price auction. The two DSPs differ in their access to consumer data: one is vertically integrated within the ad-tech stack and consequently observes consumer preferences with positive probability, while the other receives only a noisy signal. We characterize the unique Bayesian Nash equilibrium of the auction and show that greater data access for the integrated DSP softens competition, increasing expected payoffs for both bidders. At the same time, enhanced data availability improves ad relevance for consumers. We then study the effects of a mandated data-sharing remedy, deriving conditions under which forced data sharing imposes a net cost on the integrated platform even under the most favourable monetisation terms, offering a formal explanation for the persistence of data silos in digital advertising markets. \\[10pt]
    \textbf{Screening with Aggregate Quality Spillovers} \\[3pt]
    I study optimal screening on a platform where seller quality generates aggregate demand externalities. Heterogeneous sellers invest in quality at a private cost, while consumers make a participation decision based on aggregate platform quality. An exogenous bargaining parameter governs the surplus division between sellers and consumers. The feedback between quality provision and consumer participation creates a novel implementation challenge: the same nonlinear tariff can sustain multiple equilibria of the participation equation. I construct an insulating tariff that restores unique implementation. Comparing regimes in which the platform can and cannot charge a consumer fee, I show that the fee decouples the optimal quality schedule from the bargaining parameter governing surplus division. An isoelastic-uniform example yields closed-form solutions confirming that, without the fee, the platform delivers weakly lower quality than under the two-sided benchmark, with equality only at the unique bargaining split for which the optimal consumer fee is zero. The insulating construction provides a mechanism-design rationale for state-contingent platform payments such as guaranteed creator advances and developer subsidies during platform launch. \\[10pt]
    \noindent \textbf{Differential Progress on Cooperative AI} \\[3pt]
    With \href{https://lewishammond.com/}{Lewis Hammond} (Oxford/CAIF) and \href{https://www.linkedin.com/in/joss-oliver-3a49a42b1/}{Joss Oliver} (CAIF) \\[3pt]
    %This project develops a formal framework for evaluating when AI capabilities are differentially cooperative. We model capabilities as modules that transform baseline games by extending players' action spaces. Ostensibly cooperative capabilities (such as communication, commitment, simulation) exhibit dual-use potential. We posit conjectures about which standard cooperative capabilities might be differentially cooperative across different game classes, with the aim of characterizing formal conditions distinguishing capabilities that robustly improve cooperation from those whose effects depend on the baseline game.   